\chapter{Decisões: nativo, threads e memória}
\label{cap:decisoes-quarkus}

\proximasaulas

Este capítulo trata de quatro decisões que todo serviço Quarkus
enfrenta: modo JVM ou executável nativo, como conviver com reflexão,
virtual threads ou reativo, e quanto de memória dar ao container. Cada
uma aparece no mesmo formato: \emph{quando escolher X}, \emph{quando
escolher Y} e \emph{o que medir} antes de decidir. A base é a Parte~II:
quem sabe o que o JIT, o GC e o class loading fazem sabe o que se perde
e o que se ganha ao trocar cada um.

Números de desempenho aqui são do fornecedor, com data, versão e
máquina. Eles servem para dar ordem de grandeza, não para prever o seu
serviço: meça o seu, com o teste de carga da Parte~I.

\section{Modo JVM ou executável nativo}
\label{sec:jvm-vs-nativo}

O GraalVM Native Image compila a aplicação inteira, antes da execução,
num executável. O guia \citetitle{doc:graalvm-native-basics} dá a
restrição central: depois de construído o binário, nenhum elemento novo
pode ser adicionado em runtime, por exemplo por class loading. A
documentação chama isso de \termo{closed-world assumption}{Suposição de
mundo fechado: tudo o que pode executar precisa ser conhecido no build.
Só os elementos alcançáveis pela análise estática entram no binário.}.
Compare com a seção~\ref{sec:class-loading}: na HotSpot, uma classe pode
ser carregada a qualquer momento; no nativo, não existe class loading
dinâmico, e o \emph{image heap} já traz objetos criados durante o build.

O que muda em relação à Parte~\ref{parte:java-jvm}:
\begin{itemize}
\item O código é compilado antes da execução, não pelo JIT da HotSpot
  em runtime (seção~\ref{sec:jit}). O guia do Quarkus atribui a isso a
  diferença de throughput: em cargas longas, o JIT pode superar o
  nativo.
\item O GC padrão é outro (seção~\ref{sec:memoria-container}).
\item Reflexão, recursos e proxies precisam ser declarados no build
  (seção~\ref{sec:reflexao-nativo}).
\end{itemize}

\subsection*{O que o fornecedor mediu}

O guia \citetitle{doc:quarkus-building-native} publica a tabela abaixo,
atribuída ao benchmark \emph{tuned} do perf-lab do Quarkus de
21/04/2026: Quarkus~3.34.3, JDK~25.0.2, GraalVM~25.0.2-graalce, 4~CPUs,
\texttt{-Xmx512m}. A linha nativa leva o rótulo \enquote{Mandrel},
embora a configuração cite o GraalVM CE; a página não explica a
diferença. O guia não nomeia a aplicação. Os resultados dessa data
ficam em
\url{https://github.com/quarkusio/benchmarks/tree/main/results/spring-quarkus-perf-comparison/2026-04-21_09-26-13},
o que aponta para o repositório \citetitle{doc:quarkus-perf-comparison}:
um serviço REST completo, com JPA/Hibernate sobre PostgreSQL e, na
estratégia \emph{tuned}, métricas e OpenTelemetry. No comando de exemplo
do repositório, a aplicação fica em quatro núcleos fixos
(\texttt{--cpus-app 0,1,2,3}), separada do banco e do gerador de carga.
\textbf{Números do fornecedor.}

\small
\begin{longtable}{@{}p{4cm} p{2.4cm} p{3.6cm} p{1.8cm}@{}}
\toprule
Modo & 1ª requisição & Throughput de pico & RSS \\
\midrule \endhead
JVM (\emph{fast-jar}) & 4417~ms & 13\,265~tps (referência) & 304~MiB \\
JVM + cache AOT (Leyden) & 1859~ms & 12\,389~tps ($\approx$7\% abaixo) &
  240~MiB \\
Nativo (Mandrel) & 581~ms & 5\,411~tps ($\approx$59\% abaixo) & 95~MiB \\
\bottomrule
\end{longtable}
\normalsize

A \enquote{1ª requisição} é o \emph{time to first request}: do comando
de start até a primeira resposta com sucesso, não o \enquote{started
in} do log (guia \citetitle{doc:quarkus-performance-measure}). Isso
explica por que 4,4~s não contradiz o \enquote{sub-second startup} que o
guia \citetitle{doc:quarkus-building-native} atribui ao modo JVM: nessa
página, o \emph{cold start} do
modo JVM vai de cerca de 0,4~s (um REST pequeno) a cerca de 3~s (um CRUD
grande), vindo de outros posts do blog, e a medição do perf-lab inclui
subir o serviço completo e atender a primeira requisição.

Leia a tabela com a Parte~\ref{parte:java-jvm} na cabeça. O nativo ganha em tempo até a
primeira requisição e em memória; o JVM ganha em throughput de pico, que
é onde o JIT, compilando os métodos quentes com o perfil da execução
(seção~\ref{sec:jit}), faz diferença. A linha do meio usa o cache AOT
do Projeto Leyden, a mesma família de recursos dos JEPs~483 e 514
(seção~\ref{sec:class-loading}); o guia não detalha as opções usadas.
É parte do ganho de startup sem sair da HotSpot. Na mesma página, sem atribuir a fonte, o
guia estima o build nativo em 3 a 10 minutos e 4 a 8~GB de RAM na
máquina de build. Tamanhos de imagem e tempos de \emph{cold start} da
página vêm de outros posts do blog, com outras condições; por isso não
entram na tabela.

\subsection*{A decisão}

O guia resume: \enquote{Start with JVM mode. Move to native when you
have a concrete need.}

\begin{description}
\item[Quando escolher JVM] Serviço de longa duração em que o throughput
  de pico importa; build rápido no ciclo de desenvolvimento e no CI; uso
  intenso de reflexão ou class loading dinâmico, inclusive por
  bibliotecas sem extensão Quarkus.
\item[Quando escolher nativo] Startup é crítico (serverless, CLI,
  \emph{scale-to-zero}, segundo o guia); memória por réplica é o custo
  dominante; você quer a menor imagem possível.
\item[O que medir] Seguindo o guia \citetitle{doc:quarkus-performance-measure}:
  \emph{time to first request} (do comando de start até a primeira
  resposta 200, não o tempo que o framework imprime no log), RSS com
  \texttt{ps} ou \texttt{pmap} (o guia pede explicitamente para
  \emph{não} usar \texttt{docker stats}), throughput e p99 sob o teste
  de carga da Parte~\ref{parte:testes}, e o tempo de build no CI. Compare as três
  opções da tabela no \emph{seu} serviço antes de decidir.
\end{description}

\paragraph{Mandrel ou GraalVM.}
O mesmo guia descreve o Mandrel como uma distribuição \emph{downstream}
do GraalVM CE, construída sobre o OpenJDK padrão e sem suporte a
programação poliglota, e o recomenda para alvos Linux em container.
Para builds com \texttt{-Dnative} sem GraalVM instalado, há
\texttt{-Dquarkus.native.container-build=true}.

\section{Reflexão e o que quebra no nativo}
\label{sec:reflexao-nativo}

No modo JVM, reflexão custa mais que uma chamada direta; no nativo,
reflexão não declarada é um bug. A análise estática só inclui o que é
alcançável, e o que não foi alcançado não entra no binário. A página
\citetitle{doc:graalvm-metadata} explica onde está a fronteira: o Native
Image \enquote{will compute reachability at build time given that all
the call arguments are constant}, e prefere essa forma de fornecer
metadados. É o caso de \texttt{porReflexao} na seção~\ref{sec:gizmo},
que usa a constante \texttt{"privado"}. O problema real são nomes que
só existem em runtime, como os de um container que varre o classpath ao
subir. A mesma página diz ainda que, para ler ou escrever campos por
reflexão, os nomes dos campos precisam estar nos metadados; sem isso,
\texttt{Field.get} e \texttt{Field.set} lançam
\texttt{MissingReflectionRegistrationError}. O
\citetitle{doc:graalvm-compat} lista os recursos dinâmicos que exigem
metadados no build.

\subsection*{O que costuma quebrar}

O guia \citetitle{doc:quarkus-native-tips} traz os casos típicos:
\begin{description}
\item[Serialização de classes não registradas] Com Jackson, o sintoma
  citado é \texttt{No serializer found for class \ldots{} and no
  properties discovered}; com JSON-B, o JSON pode sair completamente
  vazio. O guia explica a regra geral: elementos usados via reflexão
  \enquote{are not part of the call tree so they are dead code
  eliminated}, a menos que sejam registrados.
\item[Recursos não declarados] Arquivos do classpath que devem estar no
  executável precisam ser configurados explicitamente; o caminho mais
  fácil, segundo o guia, é \texttt{quarkus.native.resources.includes}.
\item[Inicialização em build time] O guia afirma: \enquote{By default,
  Quarkus initializes all classes at build time}. Junte isso à
  seção~\ref{sec:augmentation}: o \texttt{<clinit>}
  (seção~\ref{sec:class-loading}) roda na máquina de build, e os campos
  \texttt{static} vão para o binário. Um \texttt{static final} que lê
  uma variável de ambiente guarda o valor que ela tinha na máquina de
  build. A saída é
  \texttt{--initialize-at-run-time=<pacote ou classe>}.
\item[Membros \texttt{private} em beans] Não quebram, mas custam: o guia
  \citetitle{doc:quarkus-cdi-reference} explica que acessá-los exige
  reflexão, e cada registro de reflexão aumenta o executável. Campos
  injetados \texttt{private} já são reescritos como
  \emph{package-private} por padrão (seção~\ref{sec:gizmo}).
\item[Busca programática de beans] Um bean usado só via
  \texttt{CDI.current()} pode ter sido removido no build
  (seção~\ref{sec:jandex}).
\end{description}

\subsection*{Como registrar}

\begin{itemize}
\item \texttt{@RegisterForReflection} na classe (ou numa classe de
  configuração, com \texttt{targets}, quando a classe é de uma
  biblioteca). O guia avisa que, por padrão, a anotação registra também
  classes aninhadas.
\item Um arquivo \texttt{reflect-config.json} em
  \texttt{src/main/resources/META-INF/native-image/<group-id>/<artifact-id>/},
  lido automaticamente pelo build nativo, para casos que a anotação não
  cobre; o guia avisa que, nesse caminho, manter o arquivo atualizado entre
  versões do Mandrel/GraalVM é responsabilidade sua.
\item Dentro de extensões, build items como \texttt{ReflectiveClassBuildItem}
  (guia \citetitle{doc:quarkus-writing-extensions}). É por isso que
  bibliotecas com extensão Quarkus \enquote{simplesmente funcionam} no
  nativo: alguém já escreveu esse registro.
\end{itemize}

\subsection*{A decisão}

\begin{description}
\item[Quando aceitar reflexão] Modo JVM, código de borda pouco
  executado, biblioteca madura sem alternativa. Meça antes de otimizar.
\item[Quando eliminar reflexão] Alvo nativo, caminho quente, ou quando
  existe uma extensão Quarkus para a biblioteca. Prefira injeção pelo
  construtor e campos \emph{package-private}, e escolha bibliotecas com
  extensão.
\item[O que medir] Rode a suíte de \texttt{@QuarkusIntegrationTest}
  contra o executável nativo no CI (seção~\ref{sec:testes-camadas}):
  é o teste que mostra reflexão faltando antes da produção (o
  \emph{tracing agent} do GraalVM é outro caminho, não coberto aqui). No modo JVM, o JFR e um
  profiler (seção~\ref{sec:ferramentas}) mostram se
  \texttt{Method.invoke} ou \texttt{Field.get} aparece no caminho quente.
\end{description}

\section{Virtual threads, reativo ou worker threads}
\label{sec:vt-vs-reativo}

Por baixo, o Quarkus tem um motor reativo: o guia
\citetitle{doc:quarkus-reactive-architecture} diz que ele é baseado no
Eclipse Vert.x e no Netty e cuida do I/O não bloqueante, com poucas
threads de I/O (\termo{event loop}{Laço de eventos: thread que atende
muitas conexões alternando entre elas sem bloquear; se ela bloqueia,
todas as conexões que ela atende param.}) que atendem muitas conexões
e por isso precisam ser usadas com cuidado. Sobre esse motor, o guia
\citetitle{doc:quarkus-virtual-threads} compara três modelos que podem
conviver na mesma aplicação:

\small
\begin{longtable}{@{}p{3.6cm} p{4.2cm} p{4.4cm}@{}}
\toprule
Modelo & Vantagem & Custo, segundo o guia \\
\midrule \endhead
Código síncrono em \emph{worker thread} & Código simples & A
  concorrência fica limitada pelo número de workers \\
Código reativo no event loop (\texttt{Uni}, \texttt{Multi}) & Alta
  concorrência e pouco uso de recursos & Código mais complexo \\
Código síncrono em virtual thread (\texttt{@RunOn\-Virtual\-Thread}) &
  Código simples & Risco de \emph{pinning}, monopolização e pool de
  objetos ineficiente \\
\bottomrule
\end{longtable}
\normalsize

\paragraph{O que a Parte~II já explicou.}
A seção~\ref{sec:virtual-threads} mostrou que virtual threads servem para
escala em I/O, não para velocidade; que o pinning em
\texttt{synchronized} acabou no JDK~24 (o guia do Quarkus confirma: em
Java~24+, \texttt{synchronized} \enquote{no longer cause pinning}); e
que pool via \texttt{ThreadLocal} é contraproducente. O guia acrescenta
três pontos do ponto de vista do Quarkus:
\begin{itemize}
\item \texttt{@RunOnVirtualThread} só vale em endpoints bloqueantes, por
  \texttt{@Blocking} ou pela assinatura (seção~\ref{sec:panache-rest});
  um método que devolve \texttt{Uni} já roda no event loop.
\item Monopolização: o escalonador atual não preempta tarefas, então uma
  computação longa segura a carrier; o guia avisa que, em ambiente
  restrito, isso pode levar a falta de memória e ao container ser
  encerrado. Para trabalho de CPU, o guia chama virtual threads de
  \enquote{useless and counterproductive}.
\item O ecossistema Java ainda não está todo pronto. Drivers JDBC do
  PostgreSQL anteriores ao 42.6.0 causam pinning frequente em Java~21 a
  23; em Java~24+, a maioria dos drivers JDBC não prende mais a carrier.
  Clientes Vert.x com Mutiny oferecem métodos \texttt{andAwait()} que
  bloqueiam sem prender a carrier.
\end{itemize}

\subsection*{A decisão}

\begin{description}
\item[Quando escolher virtual threads] Serviço de I/O com muitas
  requisições concorrentes que esperam banco ou HTTP, time confortável
  com código imperativo, JDK~24 ou mais novo (ou JDK~21 com as
  bibliotecas conferidas para pinning). É o caminho de menor mudança a
  partir de um serviço bloqueante.
\item[Quando escolher reativo] Streaming, muitas conexões longas
  (WebSocket, SSE), composição de várias chamadas assíncronas, ou quando
  as bibliotecas que você usa já são reativas. Exige que o time saiba não
  bloquear o event loop.
\item[Quando ficar nas worker threads] Concorrência moderada e
  previsível, código legado com \texttt{ThreadLocal} pesado, ou trabalho
  de CPU. O limite de workers funciona como proteção natural.
\item[O que medir] Sob o teste de carga da Parte~I: p99 e throughput
  nos três modelos, número de threads de plataforma e RSS. Com JFR
  (seção~\ref{sec:ferramentas}), o evento \texttt{jdk.VirtualThreadPinned}
  mostra pinning residual. Monitore também as conexões do pool do banco:
  virtual threads não criam conexões, e milhares de requisições
  concorrentes passam a disputar as mesmas poucas.
\end{description}

O repositório \citetitle{doc:quarkus-perf-comparison} tem uma variante
\texttt{quarkus3-virtual} da mesma aplicação, um bom ponto de partida
para comparar os modelos com números do seu hardware.

\section{GC e memória em container}
\label{sec:memoria-container}

\subsection*{A memória que o container vê não é o heap}

O guia \citetitle{doc:quarkus-performance-measure} mede
\termo{RSS}{\emph{Resident Set Size}: memória física de fato ocupada pelo
processo, incluindo heap e memória nativa.}, e não o heap, porque o heap
é só uma parte: a JVM também usa memória nativa para metadados de
classes, pilhas de threads, código compilado e estruturas do GC. É o RSS
que o limite do container corta. Somando o que a Parte~\ref{parte:java-jvm} e o manual do
\texttt{java} (\citetitle{doc:java-cmd}) dizem sobre cada parte:

\small
\begin{longtable}{@{}p{3cm} p{4.4cm} p{4.8cm}@{}}
\toprule
Parte & Opção & Padrão no JDK~25, segundo o manual \\
\midrule \endhead
Heap & \texttt{-XX:MaxRAMPercentage} ou \texttt{-Xmx} & 25\% da memória
  disponível (\texttt{MinRAMPercentage}, 50\%, vale para heaps pequenos,
  de cerca de 125~MB) \\
Metaspace (classes) & \texttt{-XX:MaxMetaspaceSize} & Sem limite \\
Code cache (JIT) & \texttt{-XX:ReservedCodeCacheSize} & 240~MB no
  máximo (48~MB sem compilação em camadas) \\
Pilha por thread & \texttt{-Xss} & 1024~KB em Linux/x64, 2048~KB em
  Linux/AArch64 \\
Buffers diretos & \texttt{-XX:MaxDirectMemorySize} & Escolhido pela JVM
  se não definido \\
\bottomrule
\end{longtable}
\normalsize

O mesmo manual diz que, em Linux, \texttt{UseContainerSupport} vem
ligado: a JVM detecta a memória e as CPUs do container e dimensiona os
recursos com isso. Consequência direta: num container de 1~GB sem
opções, o heap máximo fica em torno de 256~MB, e o resto sobra para
memória nativa, que pode ou não ser suficiente. E as CPUs vistas pela JVM
dimensionam os pools de threads do GC e do \texttt{ForkJoinPool};
\texttt{-XX:ActiveProcessorCount} sobrescreve esse número.

\subsection*{Modo JVM: que GC usar}

A seção~\ref{sec:gc} trouxe a regra do guia de GC da Oracle: deixe a JVM
escolher, ajuste o heap e só troque de coletor se o requisito não for
atendido. Em container, isso vira:
\begin{itemize}
\item Container pequeno, abaixo de duas CPUs ou de 1792~MB: no JDK~25
  a JVM já escolhe o Serial, o perfil que o guia associa a ele
  (\citetitle{doc:gc-ergonomics}). Quem quer o G1 nesse container
  precisa passar \texttt{-XX:+UseG1GC}, e toda medição deve registrar o
  coletor que rodou (\texttt{-Xlog:gc}, seção~\ref{sec:gc}).
\item API com requisito de p99: compare G1 e ZGC no teste de carga,
  com o mesmo limite de memória de produção.
\item Fixe \texttt{-XX:MaxRAMPercentage} (ou \texttt{-Xmx}) explicitamente
  e deixe margem para a memória nativa; ajuste olhando o RSS, não o heap.
\end{itemize}

\subsection*{Modo nativo: outro GC, outro padrão}

O guia \citetitle{doc:graalvm-memory} diz que o Serial GC é o padrão no
Native Image. O G1 existe só no Oracle GraalVM (não no Community
Edition, base do Mandrel), em Linux AMD64 e AArch64. E o detalhe que mais
surpreende: sem heap máximo definido, um executável com Serial GC usa
até 80\% do tamanho da memória física. A página fala em memória física e
não descreve como o limite do container entra nessa conta; por isso,
passe \texttt{-Xmx} em runtime (ou \texttt{-R:MaxHeapSize} no build) e
não dependa do padrão. A mesma página avisa que o RSS pode subir
temporariamente durante uma coleta, o que importa em ambientes com
limite de memória, como containers.

\subsection*{A decisão}

\begin{description}
\item[Quando dar mais memória] RSS próximo do limite sob carga, pausas de
  GC frequentes no JFR, ou heap pequeno demais para o conjunto de dados
  em uso.
\item[Quando reduzir] RSS estável muito abaixo do limite no pico de
  carga, com folga de pausa no p99. Em escala, cada megabyte por réplica
  vira custo.
\item[O que medir] RSS com \texttt{ps -o rss} ou \texttt{pmap -x} (não
  \texttt{docker stats}, segundo o guia de medição) em repouso e no pico
  do teste de carga; pausas e alocação com JFR; reinícios por
  \emph{OOMKilled} no orquestrador. Anote JDK, GC, limites de memória e
  CPU junto de cada número, como pede a seção~\ref{sec:gc}.
\end{description}

\section{Onde aprofundar}

\subsection*{Documentação}
\small
\begin{longtable}{@{}p{5cm} p{7.2cm}@{}}
\toprule
Fonte & Para quê \\
\midrule \endhead
\citetitle{doc:quarkus-building-native} & Como gerar o executável,
  Mandrel e GraalVM, e a comparação JVM/Leyden/nativo \\
\citetitle{doc:quarkus-performance-measure} & Como medir RSS e tempo até
  a primeira requisição sem se enganar \\
\citetitle{doc:graalvm-native-basics} e \citetitle{doc:graalvm-compat}
  & Mundo fechado, inicialização em build time e o que exige metadados \\
\citetitle{doc:quarkus-native-tips} & Registro de reflexão, recursos e
  inicialização em runtime \\
\citetitle{doc:graalvm-memory} & GCs e tamanho de heap no nativo \\
\citetitle{doc:quarkus-virtual-threads} e
  \citetitle{doc:quarkus-reactive-architecture} & Os três modelos de
  execução e as limitações de virtual threads \\
\citetitle{doc:java-cmd} e \citetitle{doc:gc-collectors} & Opções de
  memória e coletores da HotSpot no JDK~25 \\
\bottomrule
\end{longtable}
\normalsize

\subsection*{Ferramentas e dados}
\citetitle{doc:quarkus-perf-comparison} (a aplicação e os scripts do
benchmark) e \citetitle{doc:quarkus-benchmarks} (os resultados e
gráficos). Para medir na sua máquina: JFR, \texttt{jcmd} e
async-profiler (seção~\ref{sec:ferramentas}).

\subsection*{Pesquisa acadêmica}
Antes de comparar JVM e nativo, veja como medir. Acesso e links no
apêndice~\ref{ap:pesquisa}.
\begin{itemize}
\pesquisa{paper:statistically-rigorous-java-perf}{Intervalos de confiança
  em vez de \enquote{melhor de N}, separando startup de regime estável.}
\pesquisa{paper:vm-warmup-hot-cold}{Muitos benchmarks nunca chegam a um
  regime estável limpo, mesmo depois de aquecer.}
\pesquisa{paper:gc-myths-realities}{O tamanho do heap muda o custo total do
  GC, o que importa quando o container limita a memória.}
\pesquisa{paper:graal-partial-escape-analysis}{Por que o Graal elimina
  alocações que o C2 mantém.}
\end{itemize}

\section{Desafios práticos}

\desafio{Gere o mesmo serviço em modo JVM e nativo. Meça o tempo até a
primeira requisição e o RSS seguindo o guia de medição (não use
\texttt{docker stats}). Depois rode o teste de carga da Parte~\ref{parte:testes} por dez
minutos e compare o throughput no primeiro e no último minuto. Explique
a diferença entre os dois modos usando a seção~\ref{sec:jit}.}

\desafio{Crie uma classe com
\texttt{static final String AMBIENTE = System.getenv("AMBIENTE")} e um
endpoint que devolve o valor. Gere o executável nativo com uma variável
definida no build e rode-o com outro valor. O que o endpoint devolve?
Conserte sem tirar o campo \texttt{static}.}

\desafio{Devolva de um endpoint um objeto serializado por Jackson cuja
classe só é usada por reflexão. Encontre o erro no executável nativo e
conserte de duas formas. Qual delas você deixaria no código, e por quê?}

\desafio{Pegue um endpoint bloqueante que chama o banco e rode-o em
worker thread e com \texttt{@RunOnVirtualThread}, com 50, 500 e 5000
usuários concorrentes. Meça p99, threads de plataforma e conexões em uso
no pool. Em que ponto o gargalo deixou de ser thread e passou a ser
outra coisa?}

\desafio{Rode o serviço num container com 512~MB de limite, sem opções
de memória. Descubra o heap máximo que a JVM escolheu (\texttt{jcmd}
ou \texttt{-XX:+PrintFlagsFinal}) e compare com o RSS sob carga. Repita
com o executável nativo sem \texttt{-Xmx}. O que você passaria a
configurar explicitamente?}

